% !PRINTED 110
\weberpage{110}

Dann haben wir zunächst zu zeigen, dass
\begin{equation}
 (y+iz)^2=b+ic,
\tag{3}
\end{equation}

\noindent welche reellen Werthe auch $b,c$ haben mögen, durch reelle $y,z$ befriedigt werden kann. Die Gleichung (3) ist aber erfüllt, wenn
\begin{equation}
 y^2-z^2=b,\qquad 2yz=c,
\tag{4}
\end{equation}
also
\[
 (y^2+z^2)^2=b^2+c^2
\]
und folglich
\begin{equation}
 y^2+z^2=\sqrt{b^2+c^2}.
\tag{5}
\end{equation}
Nach IX. existirt immer eine positive Quadratwurzel $\sqrt{b^2+c^2}$. Aus (4) und (5) folgt nun
\[
 y^2=\frac{b+\sqrt{b^2+c^2}}{2},\qquad
 z^2=\frac{-b+\sqrt{b^2+c^2}}{2},
\]
und die beiden Grössen $\pm b+\sqrt{b^2+c^2}$ sind offenbar positiv, da $b^2+c^2>b^2$, also auch $\sqrt{b^2+c^2}>\sqrt{b^2}$.% [sic] Die strikten Ungleichungen gelten nicht fuer c=0: dann ist b^2+c^2=b^2 und eine der beiden Groessen gleich 0, nicht positiv. Fuer die Konstruktion genuegt Nichtnegativitaet; der Quellenfehler wird nicht stillschweigend normalisiert.
\ Wir erhalten daraus
\[
 y=\pm\sqrt{\frac{b+\sqrt{b^2+c^2}}{2}},\qquad
 z=\pm\sqrt{\frac{-b+\sqrt{b^2+c^2}}{2}},
\]
worin aber das eine Vorzeichen durch das andere bestimmt ist durch die Bedingung
\[
 2yz=c;
\]
eines der beiden Vorzeichen aber bleibt willkürlich, so dass wir nicht nur eine, sondern zwei (entgegengesetzte) Lösungen von (3) erhalten.

Es sei ferner $n$ ungerade; wir nehmen die zu lösende Gleichung in der Form an
\begin{equation}
 (y+iz)^n=b+ci.
\tag{6}
\end{equation}
Ist $c=0$, also $a$ reell, so können wir $z=0$ setzen und erhalten nach IX. einen reellen Werth von $x$. Nehmen wir aber $c$ von Null verschieden an, so folgt aus (6) (Einleitung, S. 19)
\begin{equation}
 (y-iz)^n=b-ci.
\tag{7}
\end{equation}
Multipliciren wir die beiden Gleichungen (6), (7) mit einander, so folgt
\[
 (y^2+z^2)^n=b^2+c^2,
\]
und daraus (nach IX.)
%CONT>>
